% Generated draft: requires mathematical and source review.
\documentclass[11pt]{article}
\usepackage[margin=25mm]{geometry}
\usepackage{fontspec}
\setmainfont{Latin Modern Roman}
\usepackage{amsmath,amssymb,mathtools,unicode-math}
\setmathfont{Latin Modern Math}
\usepackage{enumitem,needspace,xurl,hyperref}
\setlength{\parindent}{0pt}
\setlength{\parskip}{5pt}
\setlength{\emergencystretch}{3em}
\newcommand{\N}{\mathbb{N}}
\newcommand{\Z}{\mathbb{Z}}
\newcommand{\Q}{\mathbb{Q}}
\newcommand{\R}{\mathbb{R}}
\newcommand{\C}{\mathbb{C}}
\newcommand{\F}{\mathbb{F}}
\newcommand{\A}{\mathbb{A}}
\newcommand{\PP}{\mathbb{P}}
\newcommand{\E}{\mathbb{E}}
\newcommand{\eps}{\varepsilon}
\newcommand{\dd}{\,\mathrm{d}}
\begin{document}
\section*{UnsolvedMath ID: 30003511}
\textbf{Source problem code:} catalogue entry
\subsection{Definitions and mathematical statement}
\textbf{Stable Prasad multiplicity conjecture (source formulation).}
The supplied record concerns a \(p\)-adic Galois pair, denoted using groups \(H,G\) and a field \(F\), together with a character \(\chi\). It does not specify these data or the hypotheses they must satisfy. Let \(\phi\) be a discrete \(L\)-parameter in this intended setting, and let \(\Pi_\phi\) denote its associated stable representation, as named in the source.

Write \(\mathcal L_\chi(\phi)\) for the source's intended collection of compatible lifts \(\psi\) of \(\phi\) through the displayed map
\[
H^\vee_\chi\longrightarrow{}^LG.
\]
This notation only names the proposed indexing collection: the record does not define the domain of the parameters, the group \(H^\vee_\chi\), the map, compatibility, or the equivalence relation on lifts.

The proposed assertion is that, for every discrete parameter \(\phi\) and associated \(\Pi_\phi\) covered by the intended hypotheses,
\[
m(\Pi_\phi,\chi)
=
\frac{1}{\left|\ker^1(F;H,G)\right|}
\sum_{\psi\in\mathcal L_\chi(\phi)}
\frac{|Z(\phi)|}{|Z(\psi)|}.
\]
Here \(m(\Pi_\phi,\chi)\) is the multiplicity referred to by the source. The record does not define this multiplicity, the stable representation \(\Pi_\phi\), the object \(\ker^1(F;H,G)\), or the groups or sets \(Z(\phi)\) and \(Z(\psi)\). In particular, it does not say whether the latter notation means centralizers, component groups, or quotients of centralizers. Thus the meaning and finiteness of the cardinalities, nonvanishing of the denominators, and meaning of the sum remain unresolved. This is a transcription of the proposed identity with its missing qualifications exposed, rather than a fully specified conjecture.
\subsection{Short English statement}
The proposed stable Prasad conjecture expresses the multiplicity of a stable representation attached to a discrete \(L\)-parameter as a weighted sum over compatible parameter lifts. The supplied record omits essential definitions and hypotheses, so the formula requires source clarification and human mathematical review before it becomes a precise standalone statement.
\subsection{Sources}
\begin{itemize}
\item Local UnsolvedMath registry: supplied complete record, permanent ID \(30003511\), title \emph{Stable Prasad Multiplicity Conjecture}, source problem code \texttt{catalogue entry}, in the Oberwolfach Reports Open Problems collection. The record supplies proposed year 2017 and no named proposer. The attribution to Prasad is retained from its conjecture name; no individual authorship is independently established here.
\item Supplied report citation: \emph{Harmonic Analysis and the Trace Formula} (2018), Oberwolfach Reports, pp.\ 1551--1630, DOI \texttt{10.4171/owr/2017/25}. \url{https://doi.org/10.4171/owr/2017/25}. The supplied background describes Oberwolfach Reports as documenting mathematical workshops and their open-problem discussions.
\item Supplied literature reference: Beuzart-Plessis, ``distinction for Galois pairs.'' \url{https://arxiv.org/abs/1707.06495}.
\item Supplied literature reference: ``original OWR report.'' \url{https://doi.org/10.4171/owr/2017/25}.
\item Literature assessment dated exactly \texttt{2026-08-21}: the supplied curation labels the problem partially solved and states that Beuzart-Plessis proved important transfer-invariance and Steinberg consequences, that the stable-discrete multiplicity formula is known in selected families such as \(\mathrm{GL}_n\), \(\mathrm{U}(n)\), and \(\mathrm{GSp}_4\), and that the general conjecture remains unproved. This is supplied dated curation, not an independently verified assessment. The background explicitly requests independent verification of time-sensitive status and cited claims.
\item The registry labels the statement \texttt{corrected\_verified} and reports that the official report was checked directly and that definitions were restored in the clean statement. These are supplied provenance claims; no source was visited or independently verified in preparing this draft, and the supplied formulation still lacks essential definitions.
\end{itemize}
\end{document}
